\documentclass{article}
\usepackage{amsmath,amssymb,amsthm}
\newtheorem{theorem}{Theorem}
\newtheorem{lemma}{Lemma}
\begin{document}
\title{The Exact Constraint Budget of a Finite Cauchy Moment Fiber}
\author{}
\date{}
\maketitle

\paragraph{Status and scope.}
These are complete written elementary interpolation proofs, with a small
exact rational checker. They are not Lean-checked. They audit a proposed
moment-matching gluing mechanism, not the original HeadComplexity model.
They establish neither of the requested global asymptotic bounds and
change no model definition.

For a finite signed atomic measure $\mu$, write
\[
 m_k(\mu)=\int x^k\,d\mu(x),\qquad
 S_\mu(t)=\int\frac{d\mu(x)}{t-x}.
\]
The sample points below are distinct and outside the support interval.
Exterior points to its right suffice for every application under discussion.

\begin{lemma}[Mixed interpolation basis]
Let $p,s\geq0$, let $t_1,\ldots,t_s$ be distinct, and put
$Q(x)=\prod_{j=1}^s(t_j-x)$. The functions
\[
 1,x,\ldots,x^{p-1},\quad
 (t_1-x)^{-1},\ldots,(t_s-x)^{-1}
\]
span exactly the functions $P(x)/Q(x)$ with $\deg P<p+s$.
Empty lists and the empty product are allowed.
\end{lemma}
\begin{proof}
The polynomial part multiplied by $Q$ has degrees from $s$ through
$p+s-1$ with a nonsingular triangular leading-coefficient matrix.
The $s$ polynomials $Q(x)/(t_j-x)$ form a basis of degrees below $s$:
evaluation at $t_i$ gives a diagonal matrix with nonzero diagonal.
Thus all $p+s$ displayed numerator polynomials form a basis.
\end{proof}

\begin{theorem}[Identifiability at the atom budget]
Let each of $\mu,\nu$ have at most $a$ atoms in an interval $I$.
If $p+s\geq2a$, and
\[
 m_k(\mu)=m_k(\nu)\quad(0\leq k<p),\qquad
 S_\mu(t_j)=S_\nu(t_j)\quad(1\leq j\leq s),
\]
then $\mu=\nu$ as signed measures. Positivity of weights is not needed.
\end{theorem}
\begin{proof}
Write $\mu-\nu=\sum_{i=1}^N c_i\delta_{x_i}$ with distinct $x_i$,
where $N\leq2a$. By the lemma,
$\sum_i c_i P(x_i)/Q(x_i)=0$ for every polynomial of degree below $p+s$.
For each support point use its Lagrange polynomial of degree $N-1$,
which is one there and zero at the other support points.
This gives $c_i/Q(x_i)=0$; the denominator is nonzero.
\end{proof}

\begin{theorem}[Local freedom and its exhaustion]
Consider the $2a$ coordinates of an ordered $a$-atomic measure
$\mu=\sum_{i=1}^a w_i\delta_{x_i}$, with distinct interior nodes and
nonzero weights. For any $p+s=2a$, the map consisting of the first $p$
moments and $s$ exterior samples has nonsingular Jacobian.
Consequently, near any positive such measure, fixing the first $a$
moments leaves an $a$-dimensional smooth fiber, and its $a$ exterior
sample values are local coordinates on that fiber.
\end{theorem}
\begin{proof}
A tangent vector in the kernel would satisfy
\[
 \sum_i\bigl(\dot w_i f(x_i)+w_i\dot x_i f'(x_i)\bigr)=0
\]
for every function in the mixed basis, hence for every $f=P/Q$ with
$\deg P<2a$. Such functions can have arbitrary values and first
derivatives at the $a$ nodes. Indeed, prescribe the corresponding jets
$P(x_i)=Q(x_i)f(x_i)$ and
$P'(x_i)=Q'(x_i)f(x_i)+Q(x_i)f'(x_i)$, then apply elementary Hermite
interpolation. Choosing these jets one at a time gives
$\dot w_i=0$ and $w_i\dot x_i=0$, hence the tangent vector is zero.
The inverse function theorem gives the local-coordinate assertion.
Distinctness, interior nodes, and strictly positive weights persist in
a sufficiently small neighborhood.
\end{proof}

\begin{theorem}[Positive compression and convex sample image]
Fix $p+s=2a$ and right-exterior sample points. Every finite positive
atomic measure on a compact interval has a positive measure with at
most $a$ atoms in that interval having the same first $p$ moments and
the same $s$ samples. Consequently, after any fixed values of the first
$p$ moments are prescribed, the attainable $s$-sample vectors of
positive measures with at most $a$ atoms form a convex set.
\end{theorem}
\begin{proof}
The zero measure and measures already having at most $a$ atoms need no
change. Otherwise set $d\rho=d\mu/Q$, which is positive since $Q>0$
on the interval. Let $P_a$ be the monic polynomial of degree $a$
orthogonal to all lower-degree polynomials for $\rho$. It exists
uniquely because the moment Gram matrix is positive definite: a
nonzero polynomial of degree below $a$ cannot vanish on the more than
$a$ support points.

For completeness, $P_a$ has $a$ distinct roots $\xi_i$ in the interior
of the convex hull of the support. If its number $b$ of interior
sign-change roots were below $a$, the polynomial $R$ which is the
product of those $b$ linear factors would make $P_aR$ have a constant
sign on the support interval. It is nonzero on at least one support
point, since $P_a$ has degree $a$ and the support has more than $a$
points. Its integral would be nonzero, contradicting orthogonality to
$R$. Thus $b=a$, which also proves simplicity.

Let $L_i$ be the degree-$(a-1)$ Lagrange polynomial for the roots and
put $\lambda_i=\int L_i\,d\rho$. Polynomial division by $P_a$ and
orthogonality show
\[
 \int R\,d\rho=\sum_i\lambda_iR(\xi_i)
 \qquad(\deg R<2a).
\]
Applying this to $L_i^2$ proves
$\lambda_i=\int L_i^2\,d\rho>0$. Therefore
$\widehat\mu=\sum_i\lambda_iQ(\xi_i)\delta_{\xi_i}$ is positive.
The mixed-basis lemma proves preservation of all the stated moments
and samples. Finally, take a convex combination of two finite positive
atomic measures with the prescribed moments and apply this compression.
Its samples are the corresponding convex combination.
\end{proof}

\paragraph{Consequence for the proposed tree gluing.}
A child with $a$ positive atoms can vary its $a$ parent sample values
locally while keeping $a$ moments fixed. This is actual interpolation
freedom, but it cannot be counted twice. If two such children are then
required to agree on those $a$ samples as well as those $a$ moments,
their entire atomic measures must coincide, globally, by the preceding
identifiability theorem. A new gluing invariant must therefore use fewer
shared constraints, extra atoms, approximate matching with quantified
sign margins, or a different readout. The theorem says nothing about
the existence of such alternatives, nor about a finite sign table
whose offsets and poles may be chosen jointly.

There is therefore a useful positive statement as well: mixed moment
and sample constraints admit Gaussian compression without losing
positivity, and the sample-coordinate image of the moment fiber is
convex. Such compression does not preserve additional readouts that
depend on labels or masks of individual atoms. Applying it inside a
tree while silently retaining those descendant readouts would be a
separate, unproved step.

\end{document}
